
\begin{titlepage}
		{\centering
			\vspace{0.5cm}
			{\huge\bfseries Factor Selection and Factor Strength\\
				\Large An Application to U.S. Stock Market Return\par}
			\vspace{1cm}
			{\large A thesis submitted for the degree of\\
			Bachelor of Commerce (Honours)\\}
		\vspace{0.3cm}
	{\LARGE Version: 1.5\\}
		\vspace{1cm}
		{\small by\\}
			{\Large Zhiyuan Jiang\\\par}
\vspace{1cm}
			{\small with supervision from\par}
{\Large			Dr.~Natalia Bailey\\ Dr.~David Frazier\\}
\vfill

	\includegraphics[width = 0.3\textwidth]{plot/monash}
			
			\vfill
		{\centering Department of Econometrics and Business Statistics\\ Monash University\\Australia\\ \large \today\par}
\end{titlepage}



\chapter*{Abstract}
%	\textbf{Delete before submit: Remember what to include in the abstract: The identification of this paper, what this paper did, why we did so, what we found, and what are the meaning of those findings.}\\

In this paper, we adopt a two-steps method trying to identify risk factors for the multi-factor CAPM model.
First, we employ one new concept called factor strength to categorise factors and reduce the dimension of the candidate factor group.
And then we apply a variable selection method called Elastic net to identify factors from each smaller subgroups.
%The high dimensions of factor groups increase the computational burden when applying the variable selecting algorithm, and a common feature of high-dimension factor group is that the correlation is pervasive therefore some selecting method is no longer appropriate in that scenario.\\
When identifying risk factors for the CAPM model, two main challenges are the high-dimensionality of the factor group, and the correlation among risk factors.
Our method is designed specifically for these two challenges. 
%The high dimensionality of the candidate factor group accompany with a the correlation among factors impose extra challenges for identifying factors, and our methodology provides a potential solution aiming to these challenges.
The factors strength helps to reduce the dimensionality of the candidates' factor group and provides a guideline when evaluating the performance of our factor selecting methods.
The Elastic net technique provides a potential solution for identifying factors and reducing redundant factors when candidates are potentially correlated with each other.
In brief, we find that the Elastic net working effectively concerning reduce the redundancy in the factor groups.
This result proves the feasibility of applying Elastic net algorithm on selecting the risk factor for the CAPM model.  \\
\vspace{0.5cm}\\
 \textbf{Keywords}: CAPM, Factor Model, Factor Strength, Elastic Net, Lasso


\chapter*{Acknowledgement}
I like to acknowledge...

\chapter*{Declaration}
I declare that this research paper is my own work and has not been submitted in any form for another degree or diploma at any university or other institute of tertiary education. Information derived from the published and unpublished work of others has been acknowledged in the text and a list of references is given.\\
\vspace{0.7cm}\\
Zhiyuan Jiang

